% Composición nativa basada en Fig_9-7_Cadena_Trazabilidad.drawio. No usa recortes ni escalado de imágenes.
\clearpage
\begin{figure}[p]\centering
\begin{tikzpicture}[font=\fontsize{9.5pt}{11.5pt}\selectfont,>=stealth]
\tikzset{overview/.style={draw=lafroxFuerte,fill=lafroxFuerte!5,rounded corners=2pt,align=left,text width=8cm,minimum height=0.8cm,inner sep=5pt}}
\node[overview] (c4) at (0,0.0) {\textbf{N1} -- Origen};
\node[overview] (c6) at (0,-1.35) {\textbf{N2} -- Requisito};
\node[overview] (c9) at (0,-2.7) {\textbf{N3} -- Componente};
\node[overview] (c12) at (0,-4.050000000000001) {\textbf{N4} -- Paquete EDT};
\node[overview] (c15) at (0,-5.4) {\textbf{N5} -- Caso de prueba};
\node[overview] (c18) at (0,-6.75) {\textbf{N6} -- Cambio};
\node[overview] (c21) at (0,-8.100000000000001) {\textbf{N7} -- Despliegue};
\node[overview] (c24) at (0,-9.450000000000001) {\textbf{N8} -- Aceptación};
\draw[->,lafroxFuerte] (c4.south) -- (c6.north);
\draw[->,lafroxFuerte] (c6.south) -- (c9.north);
\draw[->,lafroxFuerte] (c9.south) -- (c12.north);
\draw[->,lafroxFuerte] (c12.south) -- (c15.north);
\draw[->,lafroxFuerte] (c15.south) -- (c18.north);
\draw[->,lafroxFuerte] (c18.south) -- (c21.north);
\draw[->,lafroxFuerte] (c21.south) -- (c24.north);
\end{tikzpicture}
\caption{Cadena de trazabilidad de extremo a extremo: vista general.}\label{fig:sd9-7}
{\fontsize{9.5pt}{12pt}\selectfont Fuente: elaboración propia a partir del Caso, sección 17.1, y de las Bases Técnicas Transversales, RT-04.04.\par}
\end{figure}\clearpage
\clearpage\begin{figure}[p]\centering
\begin{tikzpicture}[font=\fontsize{9.5pt}{12pt}\selectfont,>=stealth]
\tikzset{detail/.style={draw=lafroxFuerte,fill=lafroxFuerte!4,rounded corners=2pt,align=left,text width=5.7cm,minimum height=1.55cm,inner sep=6pt}}
\node[detail] (a0) at (0,0.0) {\textbf{N1}\\ Origen\\ Párrafo del Caso o RT\\ \textbf{Ejemplo:} Caso, cap. 18, resultado 1: retiro en < 2 h};
\node[detail] (b0) at (8.0,0.0) {\textbf{N2}\\ Requisito\\ RF · RNF · RT del T-12 (645)\\ \textbf{Ejemplo:} RF-09.01 trazabilidad hacia adelante};
\draw[->,thick,lafroxFuerte] (a0.east) -- (b0.west);
\node[detail] (a1) at (0,-3.3) {\textbf{N2}\\ Requisito\\ RF · RNF · RT del T-12 (645)\\ \textbf{Ejemplo:} RF-09.01 trazabilidad hacia adelante};
\node[detail] (b1) at (8.0,-3.3) {\textbf{N3}\\ Componente\\ M1–M12 · INT-01–15\\ \textbf{Ejemplo:} M9 Calidad y trazabilidad};
\draw[->,thick,lafroxFuerte] (a1.east) -- (b1.west);
\node[detail] (a2) at (0,-6.6) {\textbf{N3}\\ Componente\\ M1–M12 · INT-01–15\\ \textbf{Ejemplo:} M9 Calidad y trazabilidad};
\node[detail] (b2) at (8.0,-6.6) {\textbf{N4}\\ Paquete EDT\\ Formulario T-14\\ \textbf{Ejemplo:} 3.4.4 Cadena de frío, retención de lotes y retiro sanitario};
\draw[->,thick,lafroxFuerte] (a2.east) -- (b2.west);
\node[detail] (a3) at (0,-9.899999999999999) {\textbf{N4}\\ Paquete EDT\\ Formulario T-14\\ \textbf{Ejemplo:} 3.4.4 Cadena de frío, retención de lotes y retiro sanitario};
\node[detail] (b3) at (8.0,-9.899999999999999) {\textbf{N5}\\ Caso de prueba\\ CP-«ID»-nn\\ \textbf{Ejemplo:} CP-RF-09.01-01 a -05; 3.8.1 · 3.8.2};
\draw[->,thick,lafroxFuerte] (a3.east) -- (b3.west);
\end{tikzpicture}
\caption*{Figura 9.7 - detalle de conexiones 1 a 4.}
{\fontsize{9.5pt}{12pt}\selectfont Fuente: elaboración propia a partir del Caso, sección 17.1, y de las Bases Técnicas Transversales, RT-04.04.\par}
\end{figure}\clearpage
\clearpage\begin{figure}[p]\centering
\begin{tikzpicture}[font=\fontsize{9.5pt}{12pt}\selectfont,>=stealth]
\tikzset{detail/.style={draw=lafroxFuerte,fill=lafroxFuerte!4,rounded corners=2pt,align=left,text width=5.7cm,minimum height=1.55cm,inner sep=6pt}}
\node[detail] (a0) at (0,0.0) {\textbf{N5}\\ Caso de prueba\\ CP-«ID»-nn\\ \textbf{Ejemplo:} CP-RF-09.01-01 a -05; 3.8.1 · 3.8.2};
\node[detail] (b0) at (8.0,0.0) {\textbf{N6}\\ Cambio\\ Solicitud de fusión con el ID\\ \textbf{Ejemplo:} Fusión aprobada por un par};
\draw[->,thick,lafroxFuerte] (a0.east) -- (b0.west);
\node[detail] (a1) at (0,-3.3) {\textbf{N6}\\ Cambio\\ Solicitud de fusión con el ID\\ \textbf{Ejemplo:} Fusión aprobada por un par};
\node[detail] (b1) at (8.0,-3.3) {\textbf{N7}\\ Despliegue\\ Digest de la imagen\\ \textbf{Ejemplo:} Imagen firmada promovida a PROD};
\draw[->,thick,lafroxFuerte] (a1.east) -- (b1.west);
\node[detail] (a2) at (0,-6.6) {\textbf{N7}\\ Despliegue\\ Digest de la imagen\\ \textbf{Ejemplo:} Imagen firmada promovida a PROD};
\node[detail] (b2) at (8.0,-6.6) {\textbf{N8}\\ Aceptación\\ Acta T-17\\ \textbf{Ejemplo:} Resultado 1 del Caso (R18-01): simulacro de retiro, H7};
\draw[->,thick,lafroxFuerte] (a2.east) -- (b2.west);
\end{tikzpicture}
\caption*{Figura 9.7 - detalle de conexiones 5 a 7.}
{\fontsize{9.5pt}{12pt}\selectfont Fuente: elaboración propia a partir del Caso, sección 17.1, y de las Bases Técnicas Transversales, RT-04.04.\par}
\end{figure}\clearpage
\noindent\textbf{Lectura complementaria de la Figura 9.7.} Cada tarjeta desarrolla el ejemplo de un requisito crítico. La matriz del Anexo 9.C se genera desde el repositorio y el historial del pipeline.\par
